\section{存储系统}

\subsection{存储器概述与分类}

\begin{mydef}{存储器的分类维度}{mem-classify}
从多个维度对存储器进行分类：
\begin{enumerate}
    \item \textbf{按存储介质}：半导体存储器（SRAM/DRAM/ROM/Flash）、磁表面存储器（磁盘）、光存储器（光盘）。
    \item \textbf{按存取方式}：
    \begin{itemize}
        \item \textbf{随机存取（RAM）}：存取时间与物理位置无关，访问任意单元时间相等。如 SRAM、DRAM。
        \item \textbf{顺序存取（SAM）}：按物理顺序线性访问，存取时间与位置相关。如磁带。
        \item \textbf{直接存取（DAM）}：先大致定位到某个区域，再在该区域内顺序查找——兼有随机和顺序的特点。如磁盘（磁头先寻道，再在磁道上顺序旋转）。
        \item \textbf{相联存取（Associative）}：按内容而非地址访问，给定数据的一部分匹配存储单元。如 Cache 中的 TAG 比较、TLB。
    \end{itemize}
    \item \textbf{按信息的可保存性}：
    \begin{itemize}
        \item \textbf{易失性}（Volatile）：断电后信息丢失，如 SRAM、DRAM。
        \item \textbf{非易失性}（Non-volatile）：断电不丢失，如 ROM、Flash、磁盘。
    \end{itemize}
    \item \textbf{按在计算机中的位置}：寄存器（CPU 内部）、主存（内存）、辅存（外存）。
\end{enumerate}
\end{mydef}

\begin{attention}
\textbf{408 常考区分：}
\begin{itemize}
    \item 「随机存取」$\neq$「RAM」。ROM 也是随机存取的（可按地址直接访问），但它是非易失性的，不属于 RAM。
    \item 磁盘是「直接存取」，不是随机存取——因为寻道和旋转延迟使不同位置的访问时间不同。
    \item 相联存储器按\textbf{内容}检索，不是按地址；Cache 的 TAG 比较即相联查找。
\end{itemize}
\end{attention}

\begin{conclusion}{半导体存储器分类}{semi-mem-classify}
\begin{tablebox}{半导体存储器分类}
\centering
\begin{tabularx}{\linewidth}{|p{3.5cm}|p{2.8cm}|X|p{2.2cm}|}
\hline
类型 & 典型器件 & 读写特性 & 易失性 \\
\hline
RAM（随机存取）& SRAM & 可读可写 & 易失 \\
                 & DRAM & 可读可写 & 易失 \\
ROM（只读存储器）& MROM & 出厂固化的只读 & 非易失 \\
                 & PROM & 用户一次编程 & 非易失 \\
                 & EPROM & 紫外线擦除后可重写 & 非易失 \\
                 & EEPROM & 电擦除可重写 & 非易失 \\
Flash & NAND/NOR Flash & 块擦除、页编程 & 非易失 \\
\hline
\end{tabularx}
\end{tablebox}

\textbf{Flash 的两种类型：}
\begin{itemize}
    \item \textbf{NOR Flash}：支持按字节随机读，读取速度快，擦除慢，适合代码存储（如固件）。
    \item \textbf{NAND Flash}：按页读写，按块擦除，密度高，适合大容量数据存储（如 SSD、U 盘）。
\end{itemize}
\end{conclusion}

\subsection{存储系统的层次结构}

\begin{mydef}{存储层次（Memory Hierarchy）}{mem-hierarchy}
现代计算机采用多级存储结构，从上到下容量递增、速度递减、每字节成本递减：

\begin{center}
\begin{tikzpicture}[node distance=1.5cm, auto]
\node[draw, rectangle, minimum width=3cm, minimum height=0.8cm, fill=gray!10] (reg) {寄存器（CPU 内部）};
\node[draw, rectangle, minimum width=3cm, minimum height=0.8cm, fill=gray!20, below of=reg] (cache) {Cache（SRAM）};
\node[draw, rectangle, minimum width=3cm, minimum height=0.8cm, fill=gray!30, below of=cache] (mem) {主存（DRAM）};
\node[draw, rectangle, minimum width=3cm, minimum height=0.8cm, fill=gray!40, below of=mem] (disk) {辅存（磁盘/SSD）};
\node[right of=reg, node distance=4cm] (speed) {\small 访问速度：最快};
\node[right of=cache, node distance=4cm] (speed2) {\small 很快};
\node[right of=mem, node distance=4cm] (speed3) {\small 较慢};
\node[right of=disk, node distance=4cm] (speed4) {\small 最慢};
\node[left of=reg, node distance=5cm] (size) {\small 容量：最小};
\node[left of=cache, node distance=5cm] (size2) {\small 较小};
\node[left of=mem, node distance=5cm] (size3) {\small 较大};
\node[left of=disk, node distance=5cm] (size4) {\small 最大};
\end{tikzpicture}
\end{center}

存储层次通过局部性在相邻层之间保存近期或常用数据：Cache 缓存主存中的块，虚拟存储系统可把主存看作外存中进程地址空间的缓存。这里的“缓存”是功能类比，并不意味着任意时刻都满足严格的物理包含关系；写回 Cache、非包含式多级 Cache 和寄存器状态都可能使简单的“上层必为下层子集”表述失效。
\end{mydef}

\begin{conclusion}{局部性原理（Principle of Locality）}{locality}
存储层次之所以有效，根本原因是程序的\textbf{局部性原理}（CS:APP \S6.2）：

\begin{enumerate}
    \item \textbf{时间局部性}（Temporal Locality）：被引用过一次的存储位置在不久的将来很可能被再次引用。
    \begin{itemize}
        \item 典型来源：循环中的指令和数据、临时变量、频繁调用的函数。
    \end{itemize}
    \item \textbf{空间局部性}（Spatial Locality）：被引用过的存储位置的\textbf{附近}位置在不久的将来很可能被引用。
    \begin{itemize}
        \item 典型来源：顺序执行的指令、数组的顺序访问、结构体字段。
    \end{itemize}
\end{enumerate}

\textbf{「步长为 $k$ 的引用模式」}的局部性：若 $k=1$（连续访问）则空间局部性好；$k$ 越大空间局部性越差；完全不规则的访问（如哈希表）局部性最差。
\end{conclusion}

\begin{attention}
\textbf{408 高频考点：}判断一段代码的局部性。关注：
\begin{itemize}
    \item 数组访问是按行还是按列——C 语言按行存储，按行访问空间局部性好，按列访问则差。
    \item 循环体内的指令具有时间局部性（重复执行）。
    \item 函数递归调用：参数和返回地址有空间局部性（栈帧连续），递归调用本身时间局部性较弱（每次调用参数不同）。
\end{itemize}
\end{attention}

\begin{conclusion}{命中率与平均访问时间}{hit-rate}
设 $t_H$ 为上层命中访问时间。为避免“缺失总时间”和“缺失惩罚”两种口径混用，本文后续统一令 $t_M$ 表示\textbf{缺失惩罚}（不含已经付出的上层命中检查时间）：
\begin{itemize}
    \item \textbf{命中率} $H$：访问上层命中的概率。
    \item \textbf{缺失率} $M = 1-H$。
    \item \textbf{平均访问时间}：$T_{\mathrm{avg}} = t_H + (1-H)\cdot t_M = t_H + M\cdot t_M$。
\end{itemize}

两级存储（Cache--主存）情况下，$t_H$ 即 Cache 命中检查时间，$t_M$ 为缺失后从主存取一整块到 Cache 并完成本次访问的额外代价。若题目把 $t_M$ 定义为“缺失访问的总时间”（已经包含命中检查），则应改用 $T_{\mathrm{avg}}=H t_H+(1-H)t_M$，不能与上式混用。
\end{conclusion}

\subsection{主存储器的基本结构}

\begin{mydef}{存储单元与寻址}{mem-cell}
主存由 $M$ 个存储单元组成，每个单元有一个唯一的地址。常见的编址方式：
\begin{itemize}
    \item \textbf{按字节编址}：每个地址对应一个字节（8位）。现代计算机普遍采用。
    \item \textbf{按字编址}：每个地址对应一个字（如 32 位或 64 位）。
\end{itemize}
若地址线有 $n$ 根，则可寻址 $2^n$ 个存储单元。

存储芯片的关键参数：
\begin{itemize}
    \item \textbf{地址线数} $a$：确定片内寻址范围，共 $2^a$ 个单元。
    \item \textbf{数据线数} $d$：每个存储单元的位数（字长）。
    \item \textbf{容量}：$2^a \times d$ 位。
\end{itemize}
\end{mydef}

\subsubsection{存储芯片的扩展}

\begin{conclusion}{位扩展}{bit-expand}
多片存储芯片的数据线并联，实现字长的扩展。

\textbf{例：}用 $16\mathrm{K}\times 4$ 位芯片构成 $16\mathrm{K}\times 8$ 位存储器。需 2 片，地址线和控制线并联，数据线分别接 D$_0$--D$_3$ 和 D$_4$--D$_7$。

位扩展后容量变大为原来的 $k$ 倍，地址线数不变。
\end{conclusion}

\begin{conclusion}{字扩展}{word-expand}
地址线增加（高位地址通过译码器产生片选信号），实现寻址范围的扩展。

\textbf{例：}用 $16\mathrm{K}\times 8$ 位芯片构成 $64\mathrm{K}\times 8$ 位存储器。需 4 片。地址线增加 2 根（高位地址经 2-4 译码器产生 4 个片选信号 $\overline{\mathrm{CS}}_0$--$\overline{\mathrm{CS}}_3$），数据线并联即可。

字扩展后容量变大为原来的 $k$ 倍，数据线数不变。
\end{conclusion}

\begin{conclusion}{字位同时扩展}{word-bit-expand}
\textbf{例：}用 $16\mathrm{K}\times 4$ 位芯片构成 $64\mathrm{K}\times 8$ 位存储器。需 $4\times 2=8$ 片。芯片每 2 片一组做位扩展（得到 $16\mathrm{K}\times 8$ 位模块），4 组之间做字扩展。

设地址线共 $A_{15}$--$A_0$（16 根，$2^{16}=64\mathrm{K}$），其中：
\begin{itemize}
    \item $A_{13}$--$A_0$（14 根）给每组芯片的片内地址（$2^{14}=16\mathrm{K}$）；
    \item $A_{15}$--$A_{14}$（2 根）经 2-4 译码器产生 4 个片选，分别选通 4 组。
\end{itemize}
\end{conclusion}

\begin{attention}
\textbf{408 常考计算：}
\begin{itemize}
    \item 给定若干芯片和所需总容量，求所需的芯片数量和地址线分配方案。
    \item 注意区分「地址线数」和「数据线数」：位扩展改变数据线数，字扩展改变地址线数。
    \item 片选信号由高位地址经译码产生，低地址给片内寻址。区分「片内地址线」和「片选地址线」。
\end{itemize}
\end{attention}

\subsection{半导体存储器：SRAM 与 DRAM}

\begin{mydef}{SRAM（静态随机存取存储器）}{sram}
SRAM 的基本存储单元是\textbf{六管 CMOS 双稳态触发器}（6-T cell）：两个交叉耦合的反相器构成锁存器，4 个晶体管做读写控制。只要持续供电，数据就稳定保持。特点：
\begin{itemize}
    \item \textbf{速度快}：静态存储单元无需刷新，通常比 DRAM 延迟低；具体延迟取决于工艺、容量和层级，题目给出参数时应以题设为准。
    \item \textbf{功耗较高}：每个 bit 需要 $6$ 个晶体管持续耗电维持状态。
    \item \textbf{密度低}：集成度远低于 DRAM，每 bit 成本高。
    \item \textbf{用途}：CPU 内部的 Cache（L1/L2/L3），要求极低延迟。
\end{itemize}
\end{mydef}

\begin{mydef}{DRAM（动态随机存取存储器）}{dram}
DRAM 的基本存储单元是\textbf{单晶体管 + 一个微小电容}（1T1C cell）：
\begin{itemize}
    \item 数据以电荷形式存储在电容上——有电荷为「1」，无电荷为「0」。
    \item 电容会缓慢漏电，因此需要\textbf{定期刷新}（refresh）以保持数据。
    \item \textbf{刷新周期}：器件规定的保持时间内必须完成对所有行的刷新；计算题通常直接给出总刷新周期 $R$。
\end{itemize}

\textbf{DRAM 的读操作是破坏性的}——读出时会消耗电容上的电荷，读后必须立即重写（读后恢复）。

\textbf{刷新方式：}
\begin{itemize}
    \item \textbf{集中刷新}：在一个刷新周期内留出一段专门时间，一次性对所有行刷新。这段期间内不能响应 CPU 访问（「死时间」），死时间 $= \text{刷新行数} \times \text{每行刷新时间}$。
    \item \textbf{分散刷新}：把每个存储周期分成正常读写和刷新两个固定阶段，逐行刷新。没有集中的长“死时间”，但存储周期被拉长，刷新阶段不能进行正常访存。
    \item \textbf{异步刷新}：将各行刷新均匀分散到题设给定的刷新周期 $R$ 内，每隔 $R/(\text{行数})$ 左右刷新一行。它避免了集中刷新的长死时间，刷新开销也不像分散刷新那样占据每个存储周期。
\end{itemize}
\end{mydef}

\begin{conclusion}{SRAM vs DRAM 对比（408 重点）}{sram-dram-compare}
\begin{tablebox}{SRAM 与 DRAM 对比}
\centering
\begin{tabularx}{\linewidth}{|p{2.1cm}|X|X|}
\hline
特性 & SRAM & DRAM \\
\hline
基本单元 & 六管触发器 & 单管 + 电容 \\
是否需要刷新 & 否 & 是（必须在规定保持时间内刷新）\\
读写速度 & 通常较快 & 通常较慢\\
集成度 & 低（6 管/bit）& 高（1 管+1 电容/bit）\\
每 bit 成本 & 高 & 低 \\
功耗（静态）& 较高 & 较低（但需刷新功耗）\\
破坏性读出 & 否 & 是 \\
典型用途 & Cache & 主存 \\
\hline
\end{tabularx}
\end{tablebox}

\textbf{408 常见辨析：}
\begin{itemize}
    \item 「SRAM 速度快、成本高」$\to$ 用于 Cache。
    \item 「DRAM 需要刷新」$\to$ 这是 DRAM 的「动态」含义。
    \item 「DRAM 破坏性读出」$\to$ 读后必须重写。
    \item 「SRAM 不需要刷新」$\to$ 但不等于非易失——断电数据仍丢失。
\end{itemize}
\end{conclusion}

\begin{attention}
\textbf{408 高频概念区分：}
\begin{itemize}
    \item \textbf{存取周期（Memory Cycle Time）}：两次独立存储器操作之间所需的最小间隔时间。DRAM 的存取周期 $>$ 存取时间，因为读后需要恢复。
    \item \textbf{存取时间（Access Time）}：从给出地址到数据出现在数据线上的时间。
    \item \textbf{带宽（Bandwidth）}：单位时间内传输的数据量（字节/秒），$\mathrm{BW} = \text{数据宽度}/\text{存取周期}$。
\end{itemize}
\end{attention}

\subsection{双口 RAM 与多模块存储器}

\begin{mydef}{双口 RAM（Dual-Port RAM）}{dual-port-ram}
双口 RAM 有两套独立的地址线、数据线和读写控制线，允许两个设备\textbf{同时}访问不同的存储单元。若两个端口同时访问\textbf{同一存储单元}：
\begin{itemize}
    \item 同时读：允许。
    \item 一读一写或同时写：冲突，需仲裁逻辑处理（写优先或读优先）。
\end{itemize}
Cache 中的多端口实现常基于双口 RAM。
\end{mydef}

\subsubsection{多模块存储器}

\begin{conclusion}{单体多字存储器}{mono-multi-word}
在单个存储周期内从同一模块读出多个字。优点是带宽高（一次取出多条指令/数据），缺点是：
\begin{itemize}
    \item 需要更宽的数据通路（总线宽度增大）。
    \item 遇到转移指令时，多取出的指令可能被废弃，效率降低。
    \item 要求访问地址连续——若指令流发生跳转则效益丧失。
\end{itemize}
\end{conclusion}

\begin{conclusion}{多体并行存储器——交叉编址}{interleaving}
将主存分成多个独立的存储体（bank），各体有自己的地址寄存器和数据寄存器，可以\textbf{重叠}操作。关键是不同体之间的\textbf{地址分配}方式。

\subsubsection*{高位交叉编址（High-Order Interleaving）}

用地址的高位选择存储体，低位作为体内地址。即连续的存储单元位于\textbf{同一存储体}中：
\[
\text{体号} = \text{高位地址},\quad \text{体内地址} = \text{低位地址}
\]

\textbf{特点：}连续地址块通常集中在同一个存储体中，因此对单个连续地址流不能像低位交叉编址那样让相邻字在各体流水访问；它更适合让不同地址区域或不同访问流分布到不同存储体并行工作。

\subsubsection*{低位交叉编址（Low-Order Interleaving）}

用地址的低位选择存储体，高位作为体内地址。即连续的存储单元分布在\textbf{不同}存储体上：
\[
\text{体号} = \text{低位地址},\quad \text{体内地址} = \text{高位地址}
\]

\textbf{特点：}连续访存时可多个存储体\textbf{流水线}工作，大幅提高带宽，是主存储器主要的并行化技术。

\subsubsection*{带宽分析（408 高频计算）}

设存储体数量为 $m$，每个存储体的存取周期为 $T$（从给出地址到数据就绪的时间），数据宽度为 $W$ 字节。

\begin{itemize}
    \item \textbf{单体单字}带宽：$\mathrm{BW} = W/T$。
    \item \textbf{$m$ 体低位交叉流水}：每间隔 $T/m$ 启动一个体，数据连续输出。流水启动后，持续带宽为：
    \[
    \mathrm{BW}_{\mathrm{interleave}} = \frac{m\cdot W}{T}\quad (\text{流水稳态})
    \]
    即带宽提升 $m$ 倍。
\end{itemize}

\begin{tablebox}{交叉编址对比}
\centering
\begin{tabular}{|l|c|c|}
\hline
特征 & 高位交叉 & 低位交叉 \\
\hline
体号由 & 高位地址 & 低位地址 \\
连续单元分布 & 同一体内 & 不同体 \\
流水线工作 & 否 & 是 \\
提高单程序带宽 & 否 & 是 \\
典型应用 & 多道程序 & 向量/流式访存 \\
\hline
\end{tabular}
\end{tablebox}
\end{conclusion}

\begin{attention}
\textbf{408 高频计算——低位交叉编址存取的时机：}设 $m=4$，存取周期 $T$，总线传输周期 $\tau$（数据线上稳定数据所需的时间），则：
\begin{itemize}
    \item 每隔 $\tau$ 可以启动一个存储体。
    \item 若 $\tau = T/4$，则在 $T$ 时间内可以连续启动 4 个体，数据连续输出，带宽 $= 4W/T$（理想 $4$ 倍）。
    \item 若 $\tau > T/4$，则流水线会有气泡，带宽提升因子 $< m$。
    \item 408 真题常考：给定 $T$、$\tau$ 和 $m$，求连续读取 $n$ 个字所需时间。
\end{itemize}

\textbf{计算公式：}启动第一个体后，每隔 $\tau$ 启动下一个体。读取 $n$ 个字的总时间为：
\[
t = T + (n-1)\times \max(\tau,\, T/m)
\]
其中 $T$ 为第一个体存取周期（流水线填充时间），后续每字以 $\max(\tau,\, T/m)$ 的间隔取出。
\end{attention}

\subsection{Cache 高速缓冲存储器}

\begin{mydef}{Cache 的基本原理}{cache-basic}
Cache 位于 CPU 和主存之间，基于\textbf{局部性原理}：将当前最可能被访问的指令和数据放在 Cache 中，CPU 绝大部分访存请求由 Cache 满足。

Cache 按\textbf{块（Block / Cache Line）}为单位与主存交换数据。设 Cache 有 $C$ 行，每行 $B$ 字节（块大小），总容量 $C\times B$ 字节。

\medskip
\begin{center}
\begin{tikzpicture}[node distance=1.2cm, auto]
\node[draw, rectangle, minimum width=2cm, minimum height=0.7cm] (cpu) {CPU};
\node[draw, rectangle, minimum width=3cm, minimum height=0.7cm, right of=cpu, node distance=3.5cm] (cache) {Cache（SRAM）};
\node[draw, rectangle, minimum width=4cm, minimum height=0.7cm, right of=cache, node distance=5cm] (mem) {主存（DRAM）};
\draw[<->, thick] (cpu) -- node[above] {字传输} (cache);
\draw[<->, thick] (cache) -- node[above] {块传输} (mem);
\end{tikzpicture}
\end{center}

CPU 与 Cache 之间以\textbf{字}为单位传输（字长 = 机器字长），Cache 与主存之间以\textbf{块}为单位传输（块大小通常 32--128 字节）。
\end{mydef}

\subsubsection{Cache 的地址映射}

设主存地址位数为 $n$，Cache 有 $C$ 行，每行（块）$B$ 字节。主存被分成 $2^n/B$ 个块，每个块与 Cache 中的一行对应。

\begin{conclusion}{直接映射（Direct-Mapped Cache）}{direct-mapped-cache}
每个主存块只能映射到 Cache 中\textbf{唯一}的一行：$\text{Cache行号} = (\text{主存块号}) \bmod C$。

\begin{center}
\begin{tikzpicture}[node distance=0.8cm, auto]
\node[draw, rectangle, minimum width=10cm, minimum height=1.2cm, fill=gray!10] (addr) {主存地址划分（按字节编址）};
\node[draw, rectangle, minimum width=3.5cm, minimum height=0.8cm, right of=addr, node distance=0cm, fill=blue!15] (tag) {标记（Tag）};
\node[draw, rectangle, minimum width=3cm, minimum height=0.8cm, right of=tag, node distance=0cm, fill=green!15] (index) {索引（Index）};
\node[draw, rectangle, minimum width=3.5cm, minimum height=0.8cm, right of=index, node distance=0cm, fill=yellow!15] (offset) {块内偏移（Offset）};
\end{tikzpicture}
\end{center}

\begin{itemize}
    \item \textbf{块内偏移}：$\log_2 B$ 位。确定数据在块内的具体字节位置。
    \item \textbf{索引（Index）}：$\log_2 C$ 位。确定块映射到 Cache 中的哪一行。
    \item \textbf{标记（Tag）}：剩余高位。与 Cache 该行的 Tag 比较，判断是否命中。
\end{itemize}

\textbf{直接映射的优缺点：}
\begin{itemize}
    \item 硬件简单，查找速度快（只需比较一个 Tag）。
    \item 冲突缺失（Conflict Miss）率高——若两个频繁访问的块映射到同一行则互相颠簸（thrashing）。
\end{itemize}
\end{conclusion}

\begin{conclusion}{全相联映射（Fully-Associative Cache）}{fully-assoc-cache}
每个主存块可以映射到 Cache 中的\textbf{任意}一行。

主存地址只分为\textbf{标记（Tag）}和\textbf{块内偏移（Offset）}两部分——没有 Index 字段。查找时需要将 Tag 与 Cache 中所有行并行比较（相联查找）。

\begin{itemize}
    \item 冲突缺失率最低（只有容量缺失和强制缺失）。
    \item 硬件代价高——每行都要配备一个比较器，大规模 Cache 中不可行。仅用于小容量 Cache（如 TLB）。
\end{itemize}
\end{conclusion}

\begin{conclusion}{组相联映射（Set-Associative Cache）}{set-assoc-cache}
折中方案：将 Cache 分成若干组（Set），每组内有 $W$ 行（$W$ 路组相联）。每个主存块映射到\textbf{唯一的一组}，但可以放入该组中的任意一行。

\[
\text{组号} = (\text{主存块号}) \bmod S,\quad \text{其中 } S \text{ 为组数}
\]

\begin{center}
\begin{tikzpicture}[node distance=0.8cm, auto]
\node[draw, rectangle, minimum width=12cm, minimum height=1.2cm, fill=gray!10] (addr) {主存地址划分（$W$ 路组相联）};
\node[draw, rectangle, minimum width=4cm, minimum height=0.8cm, right of=addr, node distance=0cm, fill=blue!15] (tag) {标记（Tag）};
\node[draw, rectangle, minimum width=3.5cm, minimum height=0.8cm, right of=tag, node distance=0cm, fill=green!15] (index) {组号（Set Index）};
\node[draw, rectangle, minimum width=4.5cm, minimum height=0.8cm, right of=index, node distance=0cm, fill=yellow!15] {块内偏移（Offset）};
\end{tikzpicture}
\end{center}

\textbf{特殊情况：}
\begin{itemize}
    \item $W=1$ 时：就是直接映射（$S=C$，每组 1 行）。
    \item $W=C$ 时（$S=1$，所有行属于同一组）：就是全相联映射。
\end{itemize}

\textbf{地址位划分：}
\begin{itemize}
    \item \textbf{块内偏移}：$\log_2 B$ 位。
    \item \textbf{组号}：$\log_2 S = \log_2(C/W)$ 位。
    \item \textbf{标记}：$n - \log_2 S - \log_2 B$ 位。
\end{itemize}
\end{conclusion}

\begin{attention}
\textbf{408 高频考点——Cache 地址位划分的计算：}
给定主存容量、Cache 容量、块大小和相联度，计算地址中 Tag / Index / Offset 各占多少位。步骤：
\begin{enumerate}
    \item 确认主存总地址位数 $N$（按字节编址时 $N=\log_2(\text{主存容量})$）。
    \item 块内偏移位数 $=\log_2(\text{块大小})$。
    \item 计算组数 $S= \text{Cache行数} / W$，组号位数 $=\log_2 S$。
    \item Tag 位数 $= N - \log_2 S - \log_2(\text{块大小})$。
\end{enumerate}
\textbf{注意区分：「Cache 总容量」和「Cache 数据区容量」——前者还包含 Tag 和有效位等额外开销。}
\end{attention}

\subsubsection{Cache 替换算法}

当需要调入新块而组内（或全相联 Cache 中）无空闲行时，需要淘汰一个已有块。

\begin{conclusion}{LRU（Least Recently Used）替换算法}{lru-algo}
淘汰\textbf{最近最少使用}的块。基于时间局部性——最近用过的块更可能再次被使用。

\textbf{硬件实现：}
\begin{itemize}
    \item 为每行维护一个计数器（或时间戳），每次命中或新装入时将计数器清零，其他行计数器加 1。计数器最大的行即为最近最少使用。
    \item $W$ 路组相联时，每行需要 $\lceil\log_2 W\rceil$ 位的计数器。
\end{itemize}

LRU 能较好利用时间局部性，常作为教材中的基准替换算法。高相联度 Cache 中实现精确 LRU 的更新开销较大，实际硬件也常使用伪 LRU 或随机替换；题目给定哪种策略就按哪种策略模拟。
\end{conclusion}

\begin{conclusion}{FIFO（First-In First-Out）替换算法}{fifo-algo}
淘汰\textbf{最早装入}的块。不反映程序的局部性，命中率通常低于 LRU。

缺点：一个很早装入但频繁使用的块会被错误淘汰。
\end{conclusion}

\begin{conclusion}{随机替换（Random）}{random-algo}
随机选择一个块淘汰。硬件最简单，不维护任何历史信息。在大容量高相联度 Cache 中，随机替换的效果与 LRU 相差不大。
\end{conclusion}

\begin{attention}
\textbf{408 常考动手题：}给定一个访问序列和 Cache 配置（容量/块大小/相联度），模拟 Cache 的访问过程，统计命中次数和命中率。需要：
\begin{enumerate}
    \item 正确划分地址（确定每个地址对应的 Tag / Index）；
    \item 按替换策略维护每行/每组的状态；
    \item 每步判断命中/缺失。
\end{enumerate}
\textbf{常见陷阱：}地址按字节编址时，需将地址除以块大小得到块号；注意区分「组号」和「Cache 行号」。
\end{attention}

\subsubsection{Cache 的写策略}

\begin{conclusion}{写命中策略}{write-hit-policy}
当 CPU 要写入的数据在 Cache 中命中时：
\begin{itemize}
    \item \textbf{写直达（Write-Through）}：同时写入 Cache 和主存。主存数据始终与 Cache 一致（简单、一致性好），但每次写都需要访问主存，速度慢。
    \item \textbf{写回（Write-Back）}：只写入 Cache，并标记该行为「脏」（dirty）。该行被替换时再写回主存。写操作快，但 Cache 和主存可能不一致。
\end{itemize}
\textbf{408 重点：}写回需要为每行维护一个「脏位」（dirty bit）；替换脏行时需要额外的写回开销；多核环境下写回策略带来缓存一致性问题。
\end{conclusion}

\begin{conclusion}{写缺失策略}{write-miss-policy}
当 CPU 要写入的数据不在 Cache 中时：
\begin{itemize}
    \item \textbf{写分配（Write-Allocate）}：先将对应块从主存调入 Cache，再写入 Cache（即「读缺失 + 写命中」）。依据空间局部性——写入的地址附近可能再次被访问。
    \item \textbf{非写分配（Write-No-Allocate / Write-Around）}：绕过 Cache，直接写入主存，不调入 Cache。
\end{itemize}

\textbf{常见搭配：}写回 + 写分配；写直达 + 非写分配。这是两种最典型的组合。
\end{conclusion}

\begin{conclusion}{写缓冲（Write Buffer）}{write-buffer}
写直达策略中，每次写都要等候主存的慢速写入周期，严重拖慢 CPU。解决方法是在 Cache 和主存之间加一个\textbf{写缓冲}（小型 FIFO 队列）：
\begin{itemize}
    \item CPU 写 Cache 的同时将数据写入写缓冲，然后立即继续执行——主存的写入由写缓冲异步完成。
    \item 写缓冲满时 CPU 才需要 stall（停顿）。
    \item 后续读操作需先检查写缓冲中是否有待写回的对应地址数据。
\end{itemize}
\end{conclusion}

\subsubsection{Cache 性能分析}

\begin{formula}{Cache 平均访问时间}{cache-avg-time}
\begin{itemize}
    \item $t_c$：Cache 命中时间（包括 Tag 比较和选择）
    \item $t_m$：Cache 缺失后访问主存并完成取块的额外惩罚（不含已经付出的 Cache 检查时间）
    \item $h$：命中率（Hit Rate）
    \item $m=1-h$：缺失率（Miss Rate）
\end{itemize}
\[
T_{\mathrm{avg}} = t_c + m\cdot t_m
\]

\textbf{两级 Cache（L1 + L2）：}
\[
T_{\mathrm{avg}} = t_{c1} + (1-h_1)\bigl(t_{c2} + (1-h_2)t_m\bigr)
\]
其中 $t_{c1}$ 为 L1 检查时间，$t_{c2}$ 为 L1 缺失后访问 L2 并完成回填所需的额外时间，$t_m$ 为 L2 也缺失后访问主存的额外惩罚；$h_1,h_2$ 分别为 L1 命中率和“在 L1 缺失条件下”的 L2 命中率。

\textbf{注意（408 易错）：}上述公式采用“各层额外时间”口径；若题目把某层时间定义为从 CPU 发起访问到该层返回的总时间，必须按题目给定口径重新列分支，不能重复相加或漏加 L1 检查时间。
\end{formula}

\begin{attention}
\textbf{408 综合题常考——Cache 容量计算：}Cache 的总容量 $=$ 数据区容量 $+$ Tag 容量 $+$ 有效位 $+$（若有写回）脏位。

设 Cache 有 $C$ 行，每行 $B$ 字节，Tag 宽度 $T$ 位：
\begin{itemize}
    \item 数据区：$C\times B$ 字节。
    \item Tag 存储：$C\times T$ 位 $= C\times T/8$ 字节。
    \item 有效位：$C\times 1$ 位 $= C/8$ 字节。
    \item 脏位（写回时）：$C\times 1$ 位。
    \item 替换状态（组相联时）：位数取决于题设采用的实现。若明确采用每行年龄计数器，可按每行 $\lceil\log_2 W\rceil$ 位计算；精确 LRU、树形伪 LRU 等实现的状态位数并不相同。
\end{itemize}
所有项求和得 Cache 总容量。408 真题多次考察这种计算——算出 SRAM 总容量再推片数。
\end{attention}

\subsection{虚拟存储器（Virtual Memory）}

\begin{mydef}{虚拟存储器的动机与基本原理}{vm-basic}
虚拟存储器将\textbf{主存}用作\textbf{外存后备内容}的「缓存」：
\begin{itemize}
    \item 每个进程通常拥有独立的\textbf{虚拟地址空间}（如 32 位虚拟地址空间为 $2^{32}$ 字节 = 4GB），其可寻址范围可以大于实际物理主存容量。
    \item 程序使用虚拟地址（Virtual Address，VA），由硬件（MMU，Memory Management Unit）在运行时转换为物理地址（Physical Address，PA）。
    \item 虚拟地址空间中只有当前\textbf{活跃}部分驻留物理主存；未驻留页的后备内容可能来自可执行文件、映射文件或交换区，并非所有页面都必须预先写入 swap。
\end{itemize}

\textbf{与 Cache 的类比：}
\begin{tablebox}{Cache 与虚拟存储器的对应关系}
\centering
\begin{tabularx}{\linewidth}{|p{2.1cm}|X|X|}
\hline
概念 & Cache & 虚拟存储器 \\
\hline
上层（快的）& Cache & 主存（DRAM） \\
下层（慢的）& 主存（DRAM）& 磁盘/SSD \\
传送粒度 & 较小的 Cache 块 & 通常大得多的页/页框\\
缺失 & Cache Miss & Page Fault \\
缺失代价 & 取决于下一层，通常远低于缺页 & 涉及 OS；若需外存 I/O，代价高出多个数量级\\
映射与查找 & 主要由硬件完成 & OS 建立页表，MMU/TLB 协助地址转换\\
替换策略 & 硬件实现的 LRU/伪 LRU/随机等 & OS 页置换算法（如 Clock）\\
写策略 & 可采用写直达或写回 & 脏页淘汰时写回，干净页可直接丢弃\\
\hline
\end{tabularx}
\end{tablebox}

从“主存是外存内容的缓存”这一层次看，虚拟存储采用按需写回语义：对驻留页的普通写操作先修改主存并置脏，通常只在脏页被淘汰或显式同步时更新后备存储，而不会把每条存储指令都直写外存。虚拟存储主要由 OS 与 MMU 协同管理，Cache 则主要由硬件管理。
\end{mydef}

\subsubsection{页式虚拟存储器}

\begin{conclusion}{页式（Paging）虚拟存储器}{paging}
将虚拟地址空间和物理地址空间都划分为固定大小的\textbf{页（Page）}，虚拟页（Virtual Page，VP）映射到物理页框（Physical Page Frame，PP）。典型页大小：4KB（$2^{12}$ 字节）。

\textbf{虚拟地址构成：}
\begin{center}
\begin{tikzpicture}[node distance=0.8cm, auto]
\node[draw, rectangle, minimum width=10cm, minimum height=1.2cm, fill=gray!10] (addr) {虚拟地址（$N$ 位）};
\node[draw, rectangle, minimum width=5cm, minimum height=0.8cm, right of=addr, node distance=0cm, fill=blue!15] (vpn) {虚拟页号（VPN）};
\node[draw, rectangle, minimum width=5cm, minimum height=0.8cm, right of=vpn, node distance=0cm, fill=yellow!15] (vpo) {页内偏移（VPO）};
\end{tikzpicture}
\end{center}

\textbf{页表（Page Table）：}OS 为每个进程维护一张页表，记录每个虚拟页对应的物理页框号（PPN）和状态信息（有效位、脏位、访问权限等）。

\[
\text{物理地址} = \text{PPN}\;\|\;\text{VPO}
\]

页表存放于主存。每次访存理论上需要\textbf{两次}访问主存：先查页表（读 PTE），再访问数据。这是页式管理的根本开销。
\end{conclusion}

\subsubsection{TLB（Translation Lookaside Buffer）}

\begin{mydef}{TLB}{tlb}
TLB 是页表的专用 Cache，缓存最近用过的页表项（PTE）。通常采用全相联或高组相联设计，命中时间极短（$<1$ 个时钟周期）。

TLB 中的每项记录：
\begin{itemize}
    \item Tag：虚拟页号（VPN）
    \item Data：物理页框号（PPN）+ 权限位 + 脏位
\end{itemize}

TLB 命中时，翻译（VA$\to$PA）在一个时钟周期内完成，不再需要访问主存中的页表。TLB 缺失时，硬件（或 OS）从主存中的页表找到对应 PTE 并装入 TLB。
\end{mydef}

\begin{conclusion}{TLB + Cache 的综合访存流程}{tlb-cache-flow}
\textbf{这是地址转换与存储层次的常见综合题。}在同时有 TLB 和 Cache 的系统中，CPU 访存的完整流程：

\begin{enumerate}
    \item CPU 发出虚拟地址 VA。
    \item \textbf{TLB 查找：}用 VPN 查找 TLB。
    \begin{itemize}
        \item TLB 命中：获得 PPN，构成物理地址 PA。
        \item TLB 缺失：访问主存中的页表获取 PTE，装入 TLB，重新执行。
    \end{itemize}
    \item \textbf{Cache 查找：}用物理地址 PA 的 Index 和 Offset 查找 Cache。
    \begin{itemize}
        \item Cache 命中：返回数据。
        \item Cache 缺失：从主存/下级 Cache 装入块。
    \end{itemize}
\end{enumerate}

\textbf{四种典型情况的时间开销：}
\begin{tablebox}{TLB + Cache 访存分类}
\centering
\begin{tabularx}{\linewidth}{|>{\centering\arraybackslash}p{2.6cm}|>{\centering\arraybackslash}p{3.5cm}|X|}
\hline
情况 & TLB & Cache \\
\hline
最佳 & 命中 & 命中 \\
TLB 缺失后命中 & 缺失（需访存页表一次）& 命中（用页表查出的 PA 再查 Cache）\\
TLB 命中后缺失 & 命中 & 缺失（需从主存调入块）\\
最差 & 缺失 & 缺失 \\
\hline
\end{tabularx}
\end{tablebox}

\textbf{注意：}若 TLB 和 Cache 都命中，则地址转换无需查主存页表，数据可由 Cache 直接返回，无需因本次取数再访问主存。
\end{conclusion}

\begin{attention}
\textbf{常见辨析——物理寻址 Cache vs 虚拟寻址 Cache：}
\begin{itemize}
    \item \textbf{物理寻址 Cache（PIPT, Physically Indexed, Physically Tagged）：}用 PA 的 Index 选组、Tag 比较。需先完成地址翻译（TLB 查完），TLB 在关键路径上。
    \item \textbf{虚拟寻址 Cache（VIVT, Virtually Indexed, Virtually Tagged）：}直接用 VA 的 Index 和 Tag，省去 TLB。难点：虚拟地址的同义/别名（同一物理地址可能对应多个虚拟地址）。
    \item \textbf{虚拟索引物理标记 Cache（VIPT, Virtually Indexed, Physically Tagged）：}用 VA 的低位（Index + Offset）选组，PA 的 Tag 比较——Index 和 Offset 位在虚拟和物理地址中相同（因为页大小对齐）。可以在 TLB 翻译的同时开始 Cache 的组选择，TLB 翻译完成后用 PPN 的高位做 Tag 比较。这是现代处理器 L1 Cache 的常见设计。
\end{itemize}
\end{attention}

\subsubsection{段式与段页式}

\begin{conclusion}{段式（Segmentation）与段页式（Segmented Paging）}{segmentation}
\textbf{段式虚拟存储器：}按程序的逻辑结构（代码段、数据段、栈段等）划分变长段，段表（Segment Table）记录段的基址（base）和界限（limit）。虚拟地址 $=$ 段号 + 段内偏移。

优点：与程序结构一致，便于共享和保护。缺点：变长分配导致外部碎片。

\textbf{段页式虚拟存储器：}先按逻辑分段，每段内再分页。虚拟地址 $=$ 段号 + 段内页号 + 页内偏移。结合了两者的优点——逻辑清晰、无外部碎片。
\end{conclusion}

\subsubsection{缺页中断（Page Fault）}

\begin{conclusion}{缺页处理流程}{page-fault}
当页表中的有效位（Valid bit）为 0（即该虚拟页不在物理主存中）时，触发缺页异常（Page Fault）：

\begin{enumerate}
    \item CPU 保存当前指令的 PC 和必要现场。
    \item 转入 OS 的缺页处理程序（ISR / Fault Handler）。
    \item OS 确认该访问属于合法虚拟页，并在相应后备存储（可执行文件、映射文件或 swap）中定位内容；按需清零页则无需从外存读取。随后在物理主存中分配一个空闲页框。
    \item 若主存已满，OS 执行\textbf{页面置换算法}（如 LRU、Clock、改进的 Clock 算法）选出一个牺牲页写回磁盘（若脏）。
    \item 从相应后备存储读入所需页，或按页面类型执行清零填充。
    \item 更新页表（VPN$\to$PPN，标记 Valid），刷新 TLB 中对应项。
    \item 返回并重新执行引发缺页的故障指令。
\end{enumerate}

若缺页需要外存 I/O，当前进程通常阻塞，调度器可让其他就绪进程使用 CPU；缺页处理延迟仍可能通过 I/O 竞争影响系统整体性能。
\end{conclusion}

\begin{attention}
\textbf{408 常考辨析——Cache 缺失 vs 缺页：}
\begin{itemize}
    \item 在 408 通常讨论的硬件管理型 Cache 中，Cache 缺失由硬件自动从下层取块并回填，通常对 OS 透明；若题目明确给出软件管理型存储结构，则应服从题设。
    \item 缺页\textbf{由 OS 软件处理}（发出异常$\to$OS 缺页处理程序$\to$磁盘 IO$\to$页表更新$\to$返回），代价巨大。
    \item 二者在存储层次中处于不同级别：Cache 缺失是第 2--3 层之间，缺页是第 3--4 层之间。
\end{itemize}
\end{attention}

\subsection{408 高频考点速查}

\begin{conclusion}{存储系统 — 408 选择题常考点}{mem-408-keypoints}
\begin{enumerate}
    \item \textbf{SRAM vs DRAM}：SRAM 六管、快、贵、Cache；DRAM 单管+电容、需刷新、主存。DRAM 破坏性读出。
    \item \textbf{ROM 系列}：MROM/PROM/EPROM/EEPROM/Flash 的可编程性和擦除方式。
    \item \textbf{Cache 地址划分}：Tag / Index / Offset 各占多少位——给定参数必会计算。
    \item \textbf{直接映射/全相联/组相联}：映射规则、冲突缺失率对比、地址划分。
    \item \textbf{LRU 替换}：手动模拟命中/缺失过程，统计命中率。
    \item \textbf{写策略组合}：写回+写分配，写直达+非写分配。
    \item \textbf{低位交叉编址}：连续访问的流水线带宽计算（$T$、$\tau$、$m$ 的关系）。
    \item \textbf{TLB + Cache + 页表的综合访存流程}：虚拟地址$\to$TLB$\to$PA$\to$Cache$\to$数据。
    \item \textbf{缺页处理}：由 OS 而非硬件处理，属于异常（fault）类。
    \item \textbf{局部性判断}：数组按行 vs 按列访问的空间局部性差异。
\end{enumerate}
\end{conclusion}

\newpage
\section{习题}

\subsection{基础过关题}

\begin{enumerate}[leftmargin=*]

    \item \textbf{（真题风格·选择题）}
    某计算机主存按字节编址，Cache 共有 16 个行，块大小为 8 字节，采用直接映射方式。
    主存地址为 $\mathtt{1234E8F8H}$ 的单元装入 Cache 后，该单元所在 Cache 行的行号（从 0 开始）是\_\_\_\_\_\_。
    \begin{enumerate}
        \item $0$
        \item $5$
        \item $15$
        \item 无法确定
    \end{enumerate}

    \ifshowsolutions\par\noindent\textbf{解析：} 主存块号为 $\lfloor\mathtt{1234E8F8H}/8\rfloor$，直接映射行号等于块号模16，也就是取地址的第3--6位。低字节 $\mathtt{F8H}=11111000_2$，相应4位全为1，故行号为15。\boxed{\textbf{答案：第3项}}\fi

    \item \textbf{（真题风格·选择题）}
    某计算机的 Cache 共有 16 块，采用 2 路组相联映射（每组 2 块），每块大小为 32 字节，按字节编址。
    主存 129 号单元所在的主存块应装入的 Cache 组号是\_\_\_\_\_\_。
    \begin{enumerate}
        \item $0$
        \item $1$
        \item $4$
        \item $6$
    \end{enumerate}

    \ifshowsolutions\par\noindent\textbf{解析：} 主存块号为 $\lfloor129/32\rfloor=4$；Cache 有 $16/2=8$ 组，故组号为 $4\bmod8=4$。\boxed{\textbf{答案：第3项}}\fi

    \item \textbf{（真题风格·选择题）}
    某计算机主存地址空间大小为 256MB，按字节编址。数据 Cache 有 8 个 Cache 行，行长为 64B，采用直接映射方式。
    主存地址中 Tag 字段的位数是\_\_\_\_\_\_。
    \begin{enumerate}
        \item $11$
        \item $17$
        \item $19$
        \item $28$
    \end{enumerate}

    \ifshowsolutions\par\noindent\textbf{解析：} 主存为 $256\text{ MB}=2^{28}$ B，地址28位；块内偏移6位，直接映射的8行需要3位行号，故 Tag 为 $28-6-3=19$ 位。\boxed{\textbf{答案：第3项}}\fi

    \item \textbf{（真题风格·选择题）}
    下列关于 Cache 映射方式的叙述中，错误的是\_\_\_\_\_\_。
    \begin{enumerate}
        \item 直接映射中，每个主存块只能映射到 Cache 的一个特定行
        \item 全相联映射中，每个主存块可以映射到 Cache 的任意一行
        \item 组相联映射中，每组包含多行，主存块映射到特定组后，可在组内任意一行存放
        \item 组相联映射的冲突缺失率高于直接映射
    \end{enumerate}

    \ifshowsolutions\par\noindent\textbf{解析：} 在容量和块大小相同的条件下，提高相联度通常减少或不增加冲突缺失；组相联的冲突缺失率不应被断言为高于直接映射。\boxed{\textbf{答案：第4项}}\fi

    \item \textbf{（真题风格·选择题）}
    下列关于 Cache 映射方式的叙述中，错误的是\_\_\_\_\_\_。
    \begin{enumerate}
        \item 直接映射方式下，Cache 利用率最高
        \item 全相联映射方式下，Cache 命中时也需要进行 Tag 比较
        \item 组相联映射方式下，每组只有一行时即为直接映射
        \item 组相联映射方式下，每组包含所有行时即为全相联映射
    \end{enumerate}

    \ifshowsolutions\par\noindent\textbf{解析：} 直接映射的硬件简单，但不能保证所谓“利用率最高”；特定地址模式反而可能频繁冲突。全相联仍需比较 Tag；一行一组退化为直接映射，一组包含全部行退化为全相联。\boxed{\textbf{答案：第1项}}\fi

    \item \textbf{（真题风格·选择题）}
    某计算机主存按字编址，由 4 个存储体构成，采用低位交叉编址。每个存储体的存取周期为 $40\,\mathrm{ns}$，
    总线传输周期为 $10\,\mathrm{ns}$。忽略发地址等额外开销，CPU 从字地址 0 开始连续读取 4 个字，需要的总时间是\_\_\_\_\_\_。
    \begin{enumerate}
        \item $40\,\mathrm{ns}$
        \item $70\,\mathrm{ns}$
        \item $80\,\mathrm{ns}$
        \item $160\,\mathrm{ns}$
    \end{enumerate}

    \ifshowsolutions\par\noindent\textbf{解析：} 4个连续字分布到4个存储体。第一个字在40 ns后就绪，随后总线每10 ns传一个字，总时间为 $40+(4-1)\times10=70$ ns。\boxed{\textbf{答案：第2项}}\fi

    \item \textbf{（填空题）}
    某计算机主存按字节编址，主存地址为 32 位。数据 Cache 容量为 32KB，块大小为 64 字节，采用 4 路组相联映射。
    则块内偏移占 \underline{\hspace{0.8cm}} 位，组号（Set Index）占 \underline{\hspace{0.8cm}} 位，Tag 占 \underline{\hspace{1.5cm}} 位。

    \ifshowsolutions\par\noindent\textbf{答案：} 6，7，19。Cache 共 $32\text{ KiB}/64\text{ B}=512$ 行，4路时有128组，故组号7位，Tag$=32-7-6=19$位。\fi

    \item \textbf{（填空题）}
    一个由 8 个存储体构成的低位交叉编址存储器，每个存储体的存取周期为 $80\,\mathrm{ns}$，总线传输周期为 $10\,\mathrm{ns}$。
    流水线方式下连续访问时，该存储器的最大带宽为 \underline{\hspace{1.5cm}} MB/s（设每个存储体每次可读写 4 字节）。

    \ifshowsolutions\par\noindent\textbf{答案：} 400 MB/s。8个存储体恰好满足 $m\tau=T$，流水稳定后每10 ns传4 B，故带宽为 $4/(10\times10^{-9})=4\times10^8$ B/s。\fi

    \item \textbf{（简答题）}
    简述时间局部性和空间局部性的含义，并各举出一个程序中的典型例子。说明为什么存储层次结构依赖于局部性原理。

    \ifshowsolutions\par\noindent\textbf{答案：} 时间局部性指刚访问过的数据或指令近期很可能再次访问，例如循环中的计数变量和循环体指令；空间局部性指访问某地址后，邻近地址近期很可能被访问，例如顺序扫描数组。Cache、TLB 和虚拟存储只保存完整地址空间的一小部分，正是利用程序访问集中在近期工作集及邻近块的性质，才可能以较小、较快的上层存储获得较高命中率。\fi

    \item \textbf{（简答题）}
    简述 Cache 的写直达（Write-Through）和写回（Write-Back）策略的区别，并说明它们各自通常搭配的写缺失（Write Miss）策略是什么，为什么这样搭配。

    \ifshowsolutions\par\noindent\textbf{答案：} 写直达在 Cache 命中写时同时更新下层存储，内存副本较新、控制简单，但写流量较大，常与非写分配搭配，避免仅为一次写入先读入整块。写回只更新 Cache 并置脏位，块被替换时才写回下层，降低写流量，通常与写分配搭配，使后续对同一块的多次写都在 Cache 中完成。这些是常见组合而非体系结构必须遵守的唯一组合，题目若指定其他策略应以题设为准。\fi

\end{enumerate}

\subsection{真题风格综合训练}

\begin{enumerate}[leftmargin=*]

    \item \textbf{（真题风格·综合题）Cache 映射与 LRU 替换。}
    某计算机按字编址，Cache 共有 4 行，块大小为 1 个字。CPU 要访问的字地址序列为：
    \[
    0,\; 1,\; 2,\; 3,\; 1,\; 2,\; 3,\; 0,\; 1,\; 2,\; 3,\; 1,\; 2,\; 3,\; 0
    \]
    设 Cache 采用 2 路组相联映射，LRU 替换策略，初始 Cache 为空。
    \begin{enumerate}
        \item[(1)] 写出每次访问时，对应的 Cache 组号。
        \item[(2)] 逐次模拟每次访问是否命中，计算 Cache 命中次数和命中率。
        \item[(3)] 若改为 4 路组相联（全相联），同样采用 LRU 替换，命中率是多少？比较两种映射方式的命中率差异，并说明原因。
    \end{enumerate}

    \ifshowsolutions\par\noindent\textbf{答案：}
    (1) 块大小为1字，组数为 $4/2=2$，组号为块号模2。因此块0、2进组0，块1、3进组1。

    (2) 前四次访问0、1、2、3均缺失，之后11次访问全部命中；两组都恰好容纳各自映射来的两个不同块，没有发生替换。命中次数11，命中率 $11/15\approx73.33\%$。

    (3) 全相联 Cache 可同时容纳4个不同块，前四次装满后其余访问同样全部命中，命中率仍为 $11/15$。提高相联度只减少潜在冲突；本序列在2路映射下本来就没有冲突替换，因此命中率不变。\fi

    \item \textbf{（真题风格·选择题）Cache 地址结构与命中分析。}
    某计算机主存地址空间为 256MB，按字节编址。数据 Cache 容量为 8KB，块大小为 16 字节，
    采用 4 路组相联映射。
    \begin{enumerate}
        \item[(1)] 求主存地址中 Tag、组号（Set Index）和块内偏移各占多少位。
        \item[(2)] 设 CPU 依次访问以下地址（十进制）：$0,\; 16,\; 32,\; 0,\; 64,\; 16,\; 128,\; 0$。
        初始 Cache 为空。若采用 LRU 替换策略，求 Cache 命中次数和命中率。
        \item[(3)] 若替换策略改为 FIFO，求命中次数。比较 LRU 和 FIFO 的结果并说明原因。
    \end{enumerate}

    \ifshowsolutions\par\noindent\textbf{答案：}
    (1) 主存地址28位。块内偏移为4位；Cache 有 $8\text{ KiB}/16\text{ B}=512$ 行，4路时有128组，组号7位；Tag为 $28-7-4=17$ 位。

    (2) 访问地址对应的块号为 $0,1,2,0,4,1,8,0$。这些块分别进入不同组（除重复块外），命中发生在第4、6、8次，共3次，命中率 $3/8=37.5\%$。

    (3) 没有任何组发生容量超过4路的替换，故 FIFO 与 LRU 都命中3次；替换策略只有在某组需要淘汰块时才会造成差异。\fi

    \item \textbf{（真题风格·综合题）TLB + Cache + 页表综合。}
    某计算机系统按字节编址，虚拟地址 32 位，物理地址 28 位，页大小 4KB。系统具有 TLB 和两级 Cache。
    \begin{enumerate}
        \item[(1)] 虚拟地址中 VPN 和 VPO 各占多少位？物理地址中 PPN 和 PPO 各占多少位？
        \item[(2)] L1 数据 Cache 为 8 路组相联，容量 32KB，块大小 64 字节。写出物理地址中 Tag、组号和块内偏移各占多少位。
        \item[(3)] 设 TLB 采用全相联映射，共 16 项。当 CPU 执行访存指令访问虚拟地址 0x12345678 时，
        若 TLB 命中且 L1 Cache 命中，描述从虚拟地址到数据返回的完整过程（依次列出各步骤和各部件的输入/输出）。
        \item[(4)] 若 TLB 命中时 L1 Cache 缺失、L2 Cache 命中，设 TLB 命中时间 1 个周期，
        L1 命中时间 2 个周期，L2 命中时间（含从 L2 调入 L1）15 个周期。假定 TLB、L1 和 L2 按上述顺序串行访问，求这次访问的总延迟。
    \end{enumerate}

    \ifshowsolutions\par\noindent\textbf{答案：}
    (1) 页大小 $4\text{ KiB}=2^{12}$ B，故 VPO/PPO 均为12位，VPN为 $32-12=20$ 位，PPN为 $28-12=16$ 位。

    (2) L1 有 $32\text{ KiB}/64\text{ B}=512$ 行，8路时有64组。块内偏移6位，组号6位，物理 Tag 为 $28-6-6=16$ 位。

    (3) 虚拟地址拆成 VPN=\texttt{0x12345} 和 VPO=\texttt{0x678}。VPN送入全相联 TLB 比较；命中后得到 PPN，与不变的 VPO 拼成物理地址。再从物理地址取组号和 Tag 访问 L1：读出相应组的8个 Tag 并比较，命中路中的块结合块内偏移选出目标字节/字返回 CPU。因 TLB 命中，不访问内存页表；因 L1 命中，也不访问 L2。

    (4) 按题设串行计时，总延迟为 $1+2+15=18$ 个周期。这里15周期已经包含 L2 命中以及回填 L1 的时间，不能重复再加一次回填。\fi

    \item \textbf{（真题风格·选择题）LRU 替换算法。}
    某 Cache 采用 4 路组相联映射，每组有 4 个 Cache 行。对于某一特定组，初始为空。
    该组依次访问的主存块号为：$5,\; 3,\; 7,\; 5,\; 2,\; 3,\; 4,\; 3,\; 5$。
    \begin{enumerate}
        \item[(1)] 采用 LRU 替换策略，逐次写出每次访问后该组中存放的主存块号，并统计命中次数。
        \item[(2)] 采用精确 LRU 时需要记录各行怎样的信息？若只考虑能够区分 $k!$ 种新旧次序的理论最小状态编码，至少需要多少位？对本题 $k=4$ 时是多少位？
        \item[(3)] 为何在实际处理器中，4 路组相联通常使用近似 LRU（如二叉树伪 LRU）而非真正的 LRU？
    \end{enumerate}

    \ifshowsolutions\par\noindent\textbf{答案：}
    (1) 以下按“LRU$\to$MRU”列出该组状态：
    \[
    \begin{array}{c|c|c}
    \text{访问}&\text{结果}&\text{访问后状态}\\ \hline
    5&M&5\\3&M&5,3\\7&M&5,3,7\\5&H&3,7,5\\
    2&M&3,7,5,2\\3&H&7,5,2,3\\4&M&5,2,3,4\\
    3&H&5,2,4,3\\5&H&2,4,3,5
    \end{array}
    \]
    共命中4次。

    (2) 精确 LRU 必须区分组内所有有效行从最旧到最新的次序。若 $k$ 行均有效，共有 $k!$ 种排列，理论最少状态位数为 $\lceil\log_2(k!)\rceil$；$k=4$ 时为 $\lceil\log_2 24\rceil=5$ 位。实际硬件若采用成对先后关系矩阵、每行年龄计数器等直接结构，可能使用多于该理论下界的位数，还要处理无效行状态。

    (3) 真 LRU 每次命中都可能更新多行之间的全序关系，比较和更新逻辑随相联度快速增长，影响面积、功耗和关键路径。树形伪 LRU 只用约 $k-1$ 位并沿树更新，成本更低，通常能取得接近 LRU 的效果。\fi

    \item \textbf{（真题风格·选择题）低位交叉编址综合。}
    某计算机主存按字编址，每字 4 字节，由 $m = 4$ 个存储体构成，采用低位交叉编址。每个存储体的存取周期 $T = 40\,\mathrm{ns}$，
    总线传输周期 $\tau = 10\,\mathrm{ns}$，每个存储体一次读写一个字。忽略发地址等额外开销，并假定改用高位交叉编址时这 16 个连续字均落在同一存储体。
    \begin{enumerate}
        \item[(1)] 求该存储器在流水线方式下的最大带宽（MB/s）。
        \item[(2)] 若 CPU 连续从字地址 0 开始读取 16 个字（共64字节），需要多少时间？
        \item[(3)] 若改用高位交叉编址，同样读取这16个连续字，需要多少时间？比较 (2) 和 (3) 的结果，说明低位交叉编址的优势。
        \item[(4)] 若总线传输周期变为 $20\,\mathrm{ns}$，而存取周期仍为 $40\,\mathrm{ns}$，最大带宽有何变化？这个变化对 $m$ 的取值有何启示？
    \end{enumerate}

    \ifshowsolutions\par\noindent\textbf{答案：}
    (1) 启动间隔为 $\max(\tau,T/m)=\max(10,10)=10$ ns，最大带宽为 $4\text{ B}/10\text{ ns}=400$ MB/s。

    (2) 第一个字在40 ns后就绪，之后每10 ns传一个字，16个字共需
    \[
    40+(16-1)\times10=190\text{ ns}.
    \]

    (3) 按题设高位交叉时16个字落在同一体内，不能重叠该体的存取周期，共需 $16\times40=640$ ns。低位交叉把连续字分散到不同体，可流水重叠各体的存取。

    (4) 新启动间隔为 $\max(20,40/4)=20$ ns，最大带宽降为 $4/20\text{ ns}=200$ MB/s。此时只需 $m\geqslant T/\tau=2$ 就能使总线成为瓶颈，继续增加存储体不会再提高这条总线上的连续传输带宽。\fi

    \item \textbf{（真题风格·选择题）虚拟存储器与 Cache 综合。}
    某计算机按字节编址，虚拟地址 32 位，物理地址 24 位，页大小 4KB。数据 Cache 采用 2 路组相联映射，
    容量 16KB，块大小 32 字节。
    \begin{enumerate}
        \item[(1)] 虚拟地址中，VPN 占多少位？VPO 占多少位？
        \item[(2)] 物理地址中，PPN 占多少位？PPO 占多少位？
        \item[(3)] Cache 中，块内偏移占多少位？组号占多少位？Tag 占多少位？
        \item[(4)] 设某进程的页表部分内容为：虚页号 0x00123 $\to$ 实页号 0x0AB，虚页号 0x00124 $\to$ 实页号 0x0CD。
        该进程访问虚拟地址 $\mathtt{0x00123456}$。若 TLB 命中，且 Cache 也命中，写出完整的地址转换流程，
        并得出最终访问的 Cache 组号和 Tag。
    \end{enumerate}

    \ifshowsolutions\par\noindent\textbf{答案：}
    (1) VPN 20位，VPO 12位。

    (2) PPN $=24-12=12$ 位，PPO 12位。

    (3) Cache 有 $16\text{ KiB}/32\text{ B}=512$ 行，2路时有256组。块内偏移5位，组号8位，Tag $=24-8-5=11$ 位。

    (4) 虚拟地址 \texttt{0x00123456} 的 VPN 为 \texttt{0x00123}、页内偏移为 \texttt{0x456}。TLB 命中得到 PPN=\texttt{0x0AB}，拼成物理地址 \texttt{0x0AB456}。对该物理地址划分可得块内偏移 $=\texttt{0x16}$，Cache 组号 $=(\texttt{0x0AB456}\gg5)\mathbin{\&}\texttt{0xFF}=\texttt{0xA2}$，Tag $=\texttt{0x0AB456}\gg13=\texttt{0x55}$。Cache 在组 \texttt{0xA2} 的两路中比较 Tag，命中 \texttt{0x55} 后按块内偏移返回数据。\fi

    \item \textbf{（综合·局部性分析）}
    分析以下两段 C 代码的局部性表现，并估算其在 Cache 中的缺失率。
    设 Cache 块大小为 64 字节，\texttt{int} 占 4 字节，$N = 1024$，数组 \texttt{a} 按行优先存储。为使估算唯一，假定 Cache 容量小于 64 KiB，且代码 B 从一列转到下一列前，上一列涉及的各行 Cache 块均已被淘汰；忽略硬件预取和冲突细节。
    \begin{lstlisting}[language=C]
// 代码A: 按行访问
int sum = 0;
for (int i = 0; i < N; i++)
    for (int j = 0; j < N; j++)
        sum += a[i][j];

// 代码B: 按列访问
int sum = 0;
for (int j = 0; j < N; j++)
    for (int i = 0; i < N; i++)
        sum += a[i][j];
    \end{lstlisting}
    \begin{enumerate}
        \item[(1)] 分别指出代码 A 和代码 B 的时间局部性和空间局部性表现，并解释原因。
        \item[(2)] 设 Cache 初始为空，容量远小于数组大小。估算代码 A 的缺失率（忽略 \texttt{sum} 和循环变量的访问）。
        \item[(3)] 同等条件下，估算代码 B 的缺失率。
        \item[(4)] 若将代码 B 在虚拟存储系统（含 TLB、页表）中运行，除了 Cache 缺失率增加外，还可能带来什么性能问题？为什么？
    \end{enumerate}

    \ifshowsolutions\par\noindent\textbf{答案：}
    (1) 代码 A 按物理连续顺序访问，空间局部性好；每个数组元素只读一次，针对数组数据本身的时间局部性不强。代码 B 相邻访问相隔一整行，即 $1024\times4=4096$ B，空间局部性很差；同一 Cache 块中的相邻列元素要隔约1024次访问才重用，在题设容量和淘汰假设下不能保留到重用时。

    (2) 一个64 B 块含16个 int。代码 A 每装入一块后连续命中其余15个元素，忽略边界对齐影响时缺失率约为 $1/16=6.25\%$。

    (3) 代码 B 每次访问跨越4096 B，且题设保证返回下一列前旧块已被淘汰，因此几乎每次都缺失，缺失率约为100\%。

    (4) 4096 B 的步长通常正好跨页，连续内层迭代会触及大量不同虚页，容易造成 TLB 容量抖动和更多页表遍历；若物理内存不能容纳工作集，还会增加缺页和磁盘换入换出。具体是否真正缺页仍取决于驻留集和物理内存，不能仅由访问顺序断言。\fi

\end{enumerate}

\subsection{拓展阅读}

\begin{itemize}
    \item \textbf{Patterson \& Hennessy《计算机组成与设计（RISC-V 版）》}第 5 章「大而快：层次化存储系统」——完整论述存储层次、Cache 的直接映射、组相联和全相联三类映射方式及性能分析，可作为本章的扩展阅读。

    \item \textbf{Bryant \& O'Hallaron《深入理解计算机系统》（CS:APP）}第 6 章「存储器层次结构」——\S6.2 局部性原理的形式化讨论与量化实验；\S6.3--\S6.4 存储层次和 Cache 设计；\S6.5 编写 Cache 友好的代码；\S6.6 主存到 Cache 的地址映射详解。\textbf{CS:APP 对局部性的量化分析和 Cache 模拟实验是理解 408 Cache 题的极佳补充。}

    \item \textbf{CS:APP 第 9 章「虚拟内存」}——\S9.3--\S9.4 页表和 TLB 的详细工作机制；\S9.5 缺页处理和内存映射；\S9.6 地址翻译综合案例（Core i7 的完整地址翻译过程，含 TLB、多级页表、多级 Cache 的协同，是 408 TLB+Cache 综合题的理想参考）。

    \item \textbf{唐朔飞《计算机组成原理》}第 4 章「存储器」——4.1 概述，4.2 主存储器（SRAM/DRAM 对比、刷新方式），4.3 高速缓冲存储器（Cache 地址映射、替换算法、写策略），4.5 辅助存储器。不同版本术语和章节号可能略有差异，应以所用教材版本为准。

    \item \textbf{延伸思考：}现代处理器的 Cache 设计远不止 408 范围——多级 Cache 的包含策略（Inclusive、Exclusive、Non-Inclusive/Non-Exclusive）、硬件预取（Stride/Stream prefetcher）、非阻塞 Cache（MSHR）等。但 408 的核心仍是「给出参数$\to$划分地址位$\to$模拟访问序列$\to$统计命中率」这一基本功。
\end{itemize}
